\section{Related Work}
\label{sec:related}

\subsection{EEG Denoising}
\label{sec:related_denoise}

\subsubsection{Traditional EEG Denoising Methods}
Before the rise of deep learning, classical signal-processing techniques
dominated EEG denoising for a long time. Filtering, the wavelet transform, and
ICA, in particular, were widely adopted in practical EEG preprocessing because
they are theoretically mature and simple to implement. Each can effectively
improve EEG quality under suitable conditions, yet each also has inherent
limitations when confronted with complex, non-stationary, multi-source noise.

Filtering is the most basic denoising approach: it removes noise components in
specific frequency bands---such as low-frequency baseline drift ($<0.5$\,Hz) and
high-frequency EMG interference ($>40$\,Hz)---using band-pass or notch
filters~\cite{zhang2021eegdenoisenet}. These methods are computationally cheap
and suitable for real-time use, but they assume that signal and noise are
strictly separable in frequency. This assumption fails when the spectra overlap;
for example, ocular artifacts (EOG) overlap with the $\delta$ band
(0.5--4\,Hz). Moreover, an inappropriate filter design can introduce phase
distortion or remove useful information, harming subsequent analysis.

The wavelet transform, owing to its good time--frequency localization, is
another common tool. Wavelet denoising decomposes the signal into sub-bands at
multiple scales and applies thresholding to high-frequency coefficients. Unlike
conventional filtering, it adapts to signal non-stationarity and is well suited
to transient artifacts such as the EOG produced by an eye blink. However, its
performance depends strongly on the choice of wavelet basis, the number of
decomposition levels, and the thresholding strategy, and it struggles with
strongly non-additive or multi-source interference (e.g., EMG superimposed on
ambient electromagnetic noise). Over-thresholding may also remove useful neural
components and damage the original signal structure.

ICA, a blind source separation (BSS) technique, is widely used for artifact
removal. It assumes that the observed EEG is a linear mixture of statistically
independent sources and separates neural signals from ocular and muscular
artifacts by maximizing the independence of the recovered components. ICA is
particularly suited to multi-channel EEG and can extract the dominant artifact
components in an unsupervised manner. Its limitations, however, are pronounced:
(i) the independence assumption does not always hold for real EEG; (ii)
identifying artifact components typically requires manual intervention or
auxiliary criteria, increasing application complexity; and (iii) ICA is
sensitive to the number of electrodes and the amount of data, with markedly
reduced separation quality when channels are few or the noise level is
high~\cite{jiao2018sparse}.

In summary, filtering, the wavelet transform, and ICA can improve EEG quality
under specific conditions, but they generally rely on prior assumptions about
the noise (e.g., frequency distribution, independence), lack adaptivity, and
cannot automatically accommodate the differing signal characteristics of
different subjects. They therefore exhibit a serious lack of cross-subject
generalization, motivating a shift toward data-driven methods with stronger
modeling capacity~\cite{chen2025uninftm,liu2025echo}.

\subsubsection{Deep-Learning-Based Denoising Methods}
Deep learning has attracted considerable attention for EEG denoising. Models
such as convolutional neural networks (CNNs), recurrent neural networks (RNNs),
and denoising autoencoders (DAEs) recover EEG under complex noise through
end-to-end feature extraction and reconstruction.

CNNs are widely adopted because of their strong local-feature-extraction
ability. Common designs model the time-domain signal with one-dimensional
convolutions (1D-CNN) or extract features from time--frequency representations
(e.g., the short-time Fourier transform or wavelet transform) with
two-dimensional convolutions
(2D-CNN)~\cite{zhang2021eegdenoisenet}. CNNs perform well at suppressing
high-frequency EMG and EOG artifacts while preserving the local structure of
the waveform. However, the limited receptive field of convolutions focuses on
local features and captures long-range temporal dependencies poorly, with
limited ability to handle non-stationarity and cross-scale dynamics.

RNNs and their variants---long short-term memory (LSTM) and gated recurrent
units (GRU)---are used to strengthen temporal modeling. Their recurrent
structure captures long-range temporal dependencies and suits artifacts with
temporal correlation, such as EOG and baseline drift. Some studies further
introduce attention mechanisms to emphasize salient
segments~\cite{chen2022openworld}, or combine CNNs and RNNs so that
front-end convolutions extract local features while a back-end recurrent module
captures global temporal dependencies. Despite their strength in modeling
complex temporal dynamics, RNNs are prone to vanishing or exploding gradients
because of their sequential processing, and their inference is relatively slow,
increasing computational cost in practice.

DAEs, an unsupervised framework, are widely used for EEG noise suppression.
During training, a DAE adds noise to the input and learns to reconstruct the
clean signal through an encoder--decoder, thereby capturing the structural
features of the signal. Beyond the standard shallow architecture, convolutional
DAEs (CDAEs) and stacked DAEs (SDAEs) have been developed to improve nonlinear
feature extraction and denoising under complex noise. DAEs are relatively stable
on small datasets, adapt well to weakly labeled or unlabeled data, and are
simple to train~\cite{chiang2019fcn}. However, most DAEs adopt the mean squared
error (MSE) as the reconstruction loss, which struggles to model non-Gaussian,
non-additive artifacts such as EOG and EMG~\cite{ghosh2019eyeblink}, and they
lack an explicit characterization of the temporal dynamics of EEG, reducing
denoising performance across time segments and across subjects.

Overall, CNN-, RNN-, and DAE-based methods achieve high recovery quality for
specific noise types and handle short-duration, high-frequency noise and local
bursts well, but they remain limited in modeling capacity and adaptivity when
the signal distribution changes drastically, the noise is complex and
multi-source, or strong cross-subject generalization is required. How to combine
multi-scale feature extraction, temporal modeling, and domain-adaptation
strategies to improve robustness in complex environments is therefore a key
direction for further research.

\subsection{Individual Variability in EEG}
\label{sec:related_subject}

Individual variability in EEG arises from differences in physiology and
acquisition conditions. Anatomical factors---such as cortical folding patterns,
skull thickness, and scalp tissue structure---directly affect the attenuation
and spread of neural activity during volume conduction, causing the amplitude
and spatial distribution of scalp EEG to differ across
subjects~\cite{jiao2018sparse}. Physiological and cognitive factors---age, sex,
health status (e.g., fatigue, medication, neurological disease), and
psychological state (e.g., attention, mood)---further alter the spectral
composition and temporal dynamics of the signal~\cite{yin2022frequency}. Minor
differences in acquisition conditions, such as electrode placement, contact
impedance, and hardware noise level, add to the inconsistency of the data
distribution across subjects~\cite{ghosh2019eyeblink}.

Because of these differences, even under identical tasks and external
conditions, the EEG of different subjects exhibits a systematic
\emph{distribution shift}. This shift directly affects data-driven modeling: a
model trained on a single subject often degrades markedly when transferred to a
new subject, which is a key bottleneck for the large-scale deployment of EEG
analysis systems. Viewed through the lens of autonomous and adaptive systems,
this cross-subject distribution shift is a form of \emph{adaptation-space drift}:
the conditions under which a learning-based system operates change across
deployment contexts, so the model must track and adapt to them rather than assume
a fixed data distribution~\cite{gheibi2024drift}.

To address this, \emph{subject-aware} modeling has recently been introduced to
characterize and adapt to the individual variability that pervades EEG
data~\cite{jiao2018sparse}. The premise is that each subject has specific signal
characteristics, artifact-contamination patterns, and background conditions; for
example, some subjects show power-spectral densities skewed toward low
frequencies while others show more prominent high-frequency
components~\cite{yin2022frequency}. Blink rate, muscle-activity level, and
baseline-drift amplitude also differ systematically across subjects. By
explicitly injecting individual-specific information into the denoising and
feature-modeling process, subject-aware modeling helps a model adapt its signal
recovery and feature extraction, in principle overcoming the limitations of a
single subject-generalized model and improving cross-subject adaptivity and
stability.

In EEG denoising, the distribution differences between subjects give rise to a
\emph{subject gap}: even when good performance is achieved on a source subject,
the model often degrades substantially when transferred to a target
subject~\cite{chen2022openworld}. This is caused by changes in the underlying
EEG characteristics, differences in artifact-contamination patterns (e.g., EMG
intensity, EOG frequency), and varying background signal-to-noise ratios.
Designing denoising methods that perceive subject characteristics and adapt to
individual differences is therefore central to improving the cross-subject
robustness and practical usability of EEG systems~\cite{jia2025bimamba}.

\subsection{Diffusion Probabilistic Models for Signal Denoising}
\label{sec:related_diffusion}

Compared with traditional EEG denoising methods such as band-pass filtering and
ICA, DDPM-based approaches offer several advantages under complex noise. First,
DDPMs enjoy stable training: their loss is defined directly on noise prediction
and requires no discriminator, avoiding the mode collapse and training
instability that commonly affect GANs~\cite{ho2020ddpm} and improving
controllability and convergence. Second, DDPMs reconstruct signals
progressively, capturing both local and global features step by step,
preserving fine details and suppressing the information loss that denoising can
introduce~\cite{zhao2023ecg}. Third, diffusion models can in principle
approximate arbitrarily complex data distributions, flexibly modeling the
highly non-stationary and heterogeneous multi-source contamination of
EEG~\cite{zeng2023dmre2i}. Finally, the DDPM framework naturally supports
conditional generation, allowing auxiliary information such as a subject
embedding to be incorporated for subject-aware
denoising~\cite{choi2021ilvr}, thereby improving generalization and recovery
quality in cross-subject settings.

By bringing the DDPM framework to EEG denoising and combining it with
subject-aware modeling, the proposed \saddpm{} aims to improve the stability and
robustness of denoising under complex noise and high inter-subject
heterogeneity~\cite{zeng2023dmre2i}. In contrast, most current deep models
(e.g., CNNs, RNNs, and DAEs) focus on local features or short-term temporal
dynamics and lack a systematic mechanism to adapt to inter-subject domain
specificity; they are typically trained and validated on single-subject or
homogeneous datasets, without fully addressing cross-subject transfer and domain
adaptation, and they exhibit large performance fluctuations in open, real-world
environments.

While generative models such as GANs have shown potential for EEG generation and
denoising~\cite{hartmann2018eeggan,an2022autodenoising,lin2020bcg}, their
training instability, mode collapse, and evaluation difficulties limit their use
in reliability-critical applications. DDPMs, by contrast, have achieved
breakthrough results in image generation and reconstruction through stable
training, progressive recovery, and high-fidelity
reconstruction~\cite{ho2020ddpm,pan2022medical,pan2023synthesis,azqadan2023microstructure,lugmayr2022repaint}. Yet research on DDPM-based EEG denoising
remains largely unexplored: existing diffusion applications concentrate on
natural signals such as images and audio, with little systematic study of EEG
data, which are multi-source, heterogeneous, and highly non-stationary. The
progressive-reconstruction and conditional-control properties of diffusion
models make them theoretically well suited to building a noise-removal process
for complex EEG, with the potential---through domain-specific information such
as subject labels or feature encodings---to achieve adaptive denoising across
subjects and heterogeneous domains.

Beyond signal processing, generative AI is increasingly viewed as a mechanism for
adaptation in autonomous and adaptive systems~\cite{li2024genai}, including the
use of generative models to empower human-activity recognition in autonomic
systems~\cite{xu2026activity}; relatedly, deep learning has been used to manage
large adaptation spaces in self-adaptive systems~\cite{weyns2022deeplearning}.
Our subject-aware diffusion model fits this trend, using a conditional generative
model to adapt the denoising process to each individual domain. These
observations motivate the domain-specific, subject-aware diffusion approach
developed in this paper.
